\chapter{When Things Go Wrong}\label{ch:git-recovery}

The most useful fact about Git is one nobody tells beginners: \emph{almost
	nothing is ever really lost}. A commit that no branch points at is still in the
object database, and Git keeps a log of every position every branch has held.
This chapter is how to get things back, and how to find the commit that broke
something.

\section{\texorpdfstring{\cmd{reflog}: the Safety Net}{reflog: the Safety Net}}\label{sec:git-reflog}

\begin{definition}[Reflog]
	A local journal of every value \texttt{HEAD} and each branch has taken --- every
	commit, checkout, merge, rebase and reset --- with a timestamp. It is private
	to your clone and is never pushed.
\end{definition}

\begin{shell}
$ git reflog
3f8a2b1 HEAD@{0}: reset: moving to HEAD~3
9c4e7d2 HEAD@{1}: commit: Add config validation
7b1e4a9 HEAD@{2}: commit: Extract parse_line
2d8c5f0 HEAD@{3}: checkout: moving from main to feature
bc91a17 HEAD@{4}: commit: First day content only
\end{shell}

\begin{moral}
	Every ``I destroyed my work with \cmd{git reset --hard}'' has the same answer:
	\cmd{git reflog}, find the hash from before the mistake, and
	\cmd{git reset --hard <that hash>}. The commits were never deleted --- only the
	branch pointer moved, and the reflog remembers where it was.
\end{moral}

\begin{keys}[0.34]
	\cmd{git reflog}                 & \texttt{HEAD}'s history                        \\
	\cmd{git reflog show main}       & One branch's history                           \\
	\cmd{git reflog -{}-date=iso}    & With real timestamps                           \\
	\cmd{git reset -{}-hard HEAD@\{2\}} & Go back to where you were two moves ago     \\
	\cmd{git switch -c rescue HEAD@\{1\}} & Safer: give that position a new branch    \\
\end{keys}

\begin{note}
	The one thing the reflog cannot recover is work that was never committed.
	\cmd{git reset --hard} with unstaged changes, or \cmd{git clean -fd}, destroys
	those permanently --- which is why \autoref{sec:git-undo} singled them out.
	Reflog entries also expire, after 90 days by default.
\end{note}

\section{Recovering Lost Commits}\label{sec:git-recover}

\begin{keys}[0.34]
	\cmd{git reflog}                  & Almost always enough                          \\
	\cmd{git fsck -{}-lost-found}     & Find unreachable objects the reflog missed    \\
	\cmd{git show <sha>}              & Check you have the right one                  \\
	\cmd{git switch -c rescue <sha>}  & Give it a name and it is safe                 \\
	\cmd{git cherry-pick <sha>}       & Take just that commit onto your branch        \\
\end{keys}

\begin{eg}[A deleted branch]
\begin{shell}
$ git branch -D feature
Deleted branch feature (was 9c4e7d2).
$ # ...five minutes later...
$ git switch -c feature 9c4e7d2      # the hash was in the message
$ git reflog | grep feature          # or find it here
\end{shell}
	Git printed the hash when it deleted the branch, precisely so this is possible.
\end{eg}

\section{Fixing the Last Commit}\label{sec:git-amend}

\begin{keys}[0.34]
	\cmd{git commit -{}-amend}              & Redo the last commit with what is staged now \\
	\cmd{git commit -{}-amend -m "..."}     & Change only the message                      \\
	\cmd{git commit -{}-amend -{}-no-edit}  & Add staged changes, keep the message         \\
	\cmd{git commit -{}-amend -{}-author="A <a@b.c>"} & Fix the author                     \\
\end{keys}

\begin{note}
	\cmd{--amend} does not edit the commit --- it makes a new one with a new hash
	and moves the branch to it (\autoref{sec:git-objects}). So it is subject to the
	golden rule: amending a pushed commit requires a force push, and anyone who
	already pulled has the old one.
\end{note}

\section{Splitting and Reordering}\label{sec:git-split}

\begin{shell}
$ git rebase -i HEAD~4
\end{shell}

Mark the commit to divide with \cmd{edit}; when the rebase stops there:

\begin{shell}
$ git reset HEAD~          # undo the commit, keep the changes unstaged
$ git add -p               # stage the first idea only
$ git commit -m "Extract parse_line"
$ git add -A               # stage the rest
$ git commit -m "Add config validation"
$ git rebase --continue
\end{shell}

\begin{moral}
	This is \cmd{git add -p} (\autoref{sec:git-add}) doing the same job after the
	fact. A commit you made in a hurry is not permanent --- until you push it,
	the history is a draft.
\end{moral}

\section{\texorpdfstring{\cmd{bisect}}{bisect}}\label{sec:git-bisect}

\begin{definition}[Bisect]
	A binary search over history. You tell Git one commit where the bug exists and
	one where it does not; Git checks out the midpoint, you test and report, and it
	halves the range each time. A thousand commits take about ten tests.
\end{definition}

\begin{shell}
$ git bisect start
$ git bisect bad                   # the current commit is broken
$ git bisect good v1.0             # this tag was fine
Bisecting: 312 revisions left to test after this (roughly 8 steps)
[9c4e7d2...] Extract parse_line

$ make && ./run-tests.sh           # test this one
$ git bisect good                  # ...it passed
Bisecting: 156 revisions left to test after this (roughly 7 steps)

$ # ...repeat...
9c4e7d2f8a1b is the first bad commit
    Author: ...
    Date:   ...
    Extract parse_line

$ git bisect reset                 # back to where you started
\end{shell}

\begin{keys}[0.34]
	\cmd{git bisect start}          & Begin                                          \\
	\cmd{git bisect bad} / \cmd{good} & Report the current commit                    \\
	\cmd{git bisect skip}           & Cannot test this one --- it does not build     \\
	\cmd{git bisect run ./test.sh}  & Automate it: exit 0 is good, non-zero is bad   \\
	\cmd{git bisect log}            & What has been decided so far                   \\
	\cmd{git bisect reset}          & Stop and return                                \\
\end{keys}

\begin{moral}
	\cmd{git bisect run} is the whole thing in one command: write a script that
	exits 0 when the program behaves, and Git finds the guilty commit unattended.
	This is the payoff for the commit discipline of \autoref{sec:git-hygiene} ---
	bisect is only as good as the guarantee that each commit builds.
\end{moral}

\section{Archaeology}\label{sec:git-blame}

\begin{keys}[0.34]
	\cmd{git blame file.c}           & Who last changed each line, and in which commit \\
	\cmd{git blame -L 40,60 file.c}  & Just those lines                                \\
	\cmd{git blame -w -C file.c}     & Ignore whitespace; follow moved code            \\
	\cmd{git log -S'fopen'}          & Commits that added or removed that string       \\
	\cmd{git log -G'regex'}          & \dots{} matching a regular expression           \\
	\cmd{git log -L 40,60:file.c}    & The history of those lines specifically         \\
	\cmd{git log -{}-follow file.c}  & Across renames                                  \\
	\cmd{git show <sha>}             & The commit blame pointed at                     \\
\end{keys}

\begin{note}
	\cmd{git blame} is badly named --- its use is not finding someone to blame but
	finding the commit message that explains why a line is the way it is. The
	workflow is \cmd{blame} to get the hash, then \cmd{git show} to read the
	reasoning. This is the moment that repays every good commit message you ever
	wrote.
\end{note}

\begin{note}
	\cmd{git log -S} (the ``pickaxe'') answers a question \cmd{blame} cannot: when
	did this string appear or disappear? If a call to \cmd{fopen} used to be
	checked and no longer is, \cmd{blame} shows only the current line, while
	\cmd{log -S'if (!f)'} finds the commit that removed the check.
\end{note}

\section{Rewriting History Safely}\label{sec:git-rewrite}

\begin{keys}[0.34]
	\cmd{git commit -{}-amend}          & The last commit                              \\
	\cmd{git rebase -i <base>}          & Any recent series                            \\
	\cmd{git filter-repo}               & Wholesale surgery across all history         \\
	\cmd{git push -{}-force-with-lease} & Publish the rewrite, safely                  \\
\end{keys}

\begin{moral}[The golden rule, again]
	Rewriting is safe on commits nobody else has, and disruptive on commits they
	do. Before rewriting a shared branch: tell the team, do it, and have everyone
	run \cmd{git fetch \&\& git reset --hard origin/main} --- not \cmd{git pull},
	which would merge the old history back in and duplicate every commit.
\end{moral}

\section{Large Files and Secrets}\label{sec:git-bigfiles}

\begin{note}[Deleting is not enough]
	Committing a password and then deleting it in the next commit removes it from
	the working tree only. The old commit still contains it, and
	\cmd{git log -p} still shows it. The same applies to a 500\,MB file: every
	clone of the repository will download it forever.
\end{note}

\begin{keys}[0.36]
	\cmd{git filter-repo -{}-path secrets.env -{}-invert-paths} & Remove a file from all history \\
	\cmd{git filter-repo -{}-replace-text rules.txt} & Redact strings throughout       \\
	\cmd{git lfs track "*.psd"}      & Large File Storage: keep a pointer, store the blob elsewhere \\
	\cmd{git count-objects -vH}      & How big is this repository, really            \\
\end{keys}

\begin{moral}
	If a credential reached a remote, treat it as compromised and \textbf{rotate
		it}. Rewriting history is worth doing afterwards, but it does not
	un-distribute what other people already cloned, and it does not un-index what a
	scanner already found. Rotate first; clean second.
\end{moral}

\begin{note}
	\cmd{git filter-repo} replaced the old \cmd{filter-branch}, which was slow and
	notoriously easy to misuse. It is a separate install
	(\cmd{pip install git-filter-repo}). Whatever you use, do it on a fresh clone
	and keep the original until you have checked the result.
\end{note}

\begin{tryit}
	In a scratch repository, make three commits, then \cmd{git reset --hard HEAD~3}
	and confirm the work has apparently vanished. Now \cmd{git reflog}, find the
	hash, and \cmd{git reset --hard <it>}. Doing this once, deliberately, is what
	turns Git from something you are careful around into something you can
	experiment in.
\end{tryit}
